\scanpage{195}
\mspage{4.}
\deleted{\unsure{en vertu de}} utilisant de plus R2, on voit que le foncteur $X\unsure{-|\mapsto}r(X)$ des \deleted{\unsure{rev.}} objets sur $P$ cts par morceaux dans les schémas sur $r(P)$ cts par morceaux est pl. fidèle. Utilisant de plus R2 et R3, on trouve qu'il est \unsure{ess surjectif|un surjectif}. C'est donc une équivalence de catégories,
\margin{lemme à dégager ; n'utilise pas l'\unsure{hyp.|hyp} galoisienne sur $P$}
% slanted
\deleted{\unsure{ici|i.e.}} Notons que comme $G$ opère sur \unsure{$P$|$B$|$R$}, il opère sur $r(P)$, \deleted{donc} et on trouve que l'équivalence est compatible avec les opérations de $G$. On trouve ainsi :

\uline{Prop} 3 \inserted{Soit $P$ objet galoisien de groupe $G$ \unsure{de|dans} $\mathscr R_S$.}
Le foncteur $X\unsure{\leadsto|\mapsto}r(X\times P)/r(P)$ induit une \uline{équivalence} de la catégorie des objets de $\mathscr R_S$ subordonnés à $P$, avec la catégorie des schémas \inserted{finis} constants par morceaux sur $r(P)$, sur lesquels $G$ opère de façon compatible avec ses opérations sur $r(P)$.

Utilisant la théorie de la descente, il vient

\uline{Th.} 1 Soit $P$ objet galoisien de groupe $G$ en $\mathscr R_S$. Le foncteur $X\longmapsto r(X\times P)/r(P)$ induit une équivalence de la catégorie des objets de $\mathscr R_S$ loc. subordonnés à $P$, avec la catégorie des \uline{revêtements étales} \inserted{\deleted{\unsure{$Z$}}$Z'$} de $Z=r(P)$, \uline{à action de} $G$ compatibles avec son action sur $r(P)$, et satisfaisant la condition suivante (\unsure{automatiqu$^{\rm t}$|automatique|automatiquement} satisfaite si $r(P)=Z$ est fini sur $S$) : les $S'\unsure{-|\to}S$ étales tels que $Z'_{S'}$ soit \deleted{\unsure{étale sur|étales sur}} constant par morceaux sur $Z_{S'}$ \uline{recouvrent} $S$. / Localisant sur $S$, on trouve de même une description \deleted{de}\inserted{du \uline{champ}} $\mathscr R^P$ en termes du schéma $Z$ à groupe d'opérateurs $G$ sur $S$.
% archive sheet number: 194
